%%%%%%%%%%%%%%%%%%%%%%%%%%%%%%%%%%%%%%%%%
% Simple Sectioned Essay Template
% LaTeX Template
%
% This template has been downloaded from:
% http://www.latextemplates.com
%
% Note:
% The \lipsum[#] commands throughout this template generate dummy text
% to fill the template out. These commands should all be removed when 
% writing essay content.
%
%%%%%%%%%%%%%%%%%%%%%%%%%%%%%%%%%%%%%%%%%

%----------------------------------------------------------------------------------------
%	PACKAGES AND OTHER DOCUMENT CONFIGURATIONS
%----------------------------------------------------------------------------------------

\documentclass[12pt]{article} % Default font size is 12pt, it can be changed here

\usepackage{geometry} % Required to change the page size to A4
\geometry{a4paper} % Set the page size to be A4 as opposed to the default US Letter

\usepackage{graphicx} % Required for including pictures

\usepackage{float} % Allows putting an [H] in \begin{figure} to specify the exact location of the figure
\usepackage{wrapfig} % Allows in-line images such as the example fish picture

\usepackage{lipsum} % Used for inserting dummy 'Lorem ipsum' text into the template

\usepackage{hyperref}

\linespread{1.2} % Line spacing

%\setlength\parindent{0pt} % Uncomment to remove all indentation from paragraphs

\graphicspath{{Pictures/}} % Specifies the directory where pictures are stored

\setcounter{secnumdepth}{4}
\setcounter{tocdepth}{4}
%======================================================================================
%      EXTRA PREAMBLE INCLUSIONS
%======================================================================================

%--------------------------------------------------------------------------------------
%       CODE SNIPPETS
%--------------------------------------------------------------------------------------

\usepackage{listings} % Required for inserting code snippets
\usepackage[usenames,dvipsnames]{color} % Required for specifying custom colors and referring to colors by name

\definecolor{DarkGreen}{rgb}{0.0,0.4,0.0} % Comment color
\definecolor{highlight}{RGB}{250,220,150} % Code highlight color

\lstdefinestyle{Style1}{ % Define a style for your code snippet, multiple definitions can be made if, for example, you wish to insert multiple code snippets using different programming languages into one document
language=C++, % Detects keywords, comments, strings, functions, etc for the language specified
backgroundcolor=\color{highlight}, % Set the background color for the snippet - useful for highlighting
basicstyle=\footnotesize\ttfamily, % The default font size and style of the code
breakatwhitespace=false, % If true, only allows line breaks at white space
breaklines=true, % Automatic line breaking (prevents code from protruding outside the box)
captionpos=b, % Sets the caption position: b for bottom; t for top
commentstyle=\usefont{T1}{pcr}{m}{sl}\color{DarkGreen}, % Style of comments within the code - dark green courier font
deletekeywords={}, % If you want to delete any keywords from the current language separate them by commas
%escapeinside={\%}, % This allows you to escape to LaTeX using the character in the bracket
firstnumber=1, % Line numbers begin at line 1
frame=single, % Frame around the code box, value can be: none, leftline, topline, bottomline, lines, single, shadowbox
frameround=, % Rounds the corners of the frame for the top left, top right, bottom left and bottom right positions
keywordstyle=\color{Blue}\bf, % Functions are bold and blue
morekeywords={}, % Add any functions no included by default here separated by commas
numbers=left, % Location of line numbers, can take the values of: none, left, right
numbersep=1.2em, % Distance of line numbers from the code box
numberstyle=\tiny\color{Gray}, % Style used for line numbers
rulecolor=\color{black}, % Frame border color
showstringspaces=false, % Don't put marks in string spaces
showtabs=false, % Display tabs in the code as lines
stepnumber=1, % The step distance between line numbers, i.e. how often will lines be numbered
stringstyle=\color{Purple}, % Strings are purple
tabsize=2, % Number of spaces per tab in the code
}

% Create a command to cleanly insert a snippet with the style above anywhere in the document
\newcommand{\insertcode}[2]{\begin{itemize}\item[]\lstinputlisting[caption=#2,label=#1,style=Style1]{#1}\end{itemize}} % The first argument is the script location/filename and the second is a caption for the listing

%%%%% EXAMPLE USE CASE %%%%%
% The first argument is the script location/filename and the second is a caption for the listing

% \insertcode{"Scripts/example.pl"}{Nena would be proud.} 

%----------------------------------------------------------------------------------------







\begin{document}

%----------------------------------------------------------------------------------------
%	TITLE PAGE
%----------------------------------------------------------------------------------------

\begin{titlepage}

\newcommand{\HRule}{\rule{\linewidth}{0.5mm}} % Defines a new command for the horizontal lines, change thickness here

\center % Center everything on the page

\textsc{\LARGE Monash University}\\[1.5cm] % Name of your university/college
\textsc{\Large Electrical Engineering}\\[0.5cm] % Major heading such as course name
\textsc{\large Final Year Project (ECE4093)}\\[0.5cm] % Minor heading such as course title

\HRule \\[0.4cm]
{ \huge \bfseries Design Document}\\[0.4cm] % Title of your document
\HRule \\[1.5cm]

\begin{minipage}{0.4\textwidth}
\begin{flushleft} \large
\emph{Author:}\\
James \textsc{Anastasiou}\\ % Your name
\emph{ID:} 23438940
\end{flushleft}
\end{minipage}
~
\begin{minipage}{0.4\textwidth}
\begin{flushright} \large
\emph{Supervisor:} \\
Dr. David \textsc{Boland} % Supervisor's Name
\\
\end{flushright}
\end{minipage}\\[4cm]

{\large \today}\\[3cm] % Date, change the \today to a set date if you want to be precise

%\includegraphics{Logo}\\[1cm] % Include a department/university logo - this will require the graphicx package

\vfill % Fill the rest of the page with whitespace

\end{titlepage}

%----------------------------------------------------------------------------------------
%	TABLE OF CONTENTS
%----------------------------------------------------------------------------------------

\tableofcontents % Include a table of contents

\newpage % Begins the essay on a new page instead of on the same page as the table of contents 

%----------------------------------------------------------------------------------------
%       CONTENT	
%----------------------------------------------------------------------------------------
%!TEX root = ./main.tex
%----------------------------------------------------------------------------------------
%	INTRODUCTION
%----------------------------------------------------------------------------------------

\section{Introduction} % Major section


%------------------------------------------------

\subsection{Purpose and Scope} % Sub-section

This document aims to provide a framework which describes the process by which my final year project will be completed, and by extension the re-evaluation of goals initially set forth in the requirements specification. In order to best achieve this aim, this document includes:

\begin{itemize}
\item Project Description
\item Re-evaluation of the inital goals
\item Project Status Overview
\item Current Position Review 
\item Timeline to Expected Completion
\end{itemize}


\pagebreak

%!TEX root = ./main.tex
%---------------------------------------------------------------------------------------- PROJECT
%DESCRIPTION
%----------------------------------------------------------------------------------------

\section{Project Description} 
My final year project aimed to create a set of tools for static analytics to help determine which
algorithms within a code-base may be easily parallelized. The specific intention of it is to provide
users with a report describing which parts of their code-base may see a potential speed-up from
redeveloping them as CUDA (GPU)\cite{cuda} kernels. In order to achieve this both benchmarking and
theoretical analysis was required, as well as static analysis of the C description of the algorithm
itself. My final project utilised a combined strategy of static AST analysis, via a hook into the
Clang compiler, as well as integration with a benchmarking suite I created that utilised real high
performance libraries for both Host and Device code. The use case of this software is for developers
working primarily in modelling and mathematically intensive areas, as these tend to provide
significant opportunity to provide improvement. The expectation of potential users is that they will
note that the report alerts them to potential speedups, and then use the performance metrics to
determine whether or not to hire a specialised engineer to redesign the relevant components of their
system.

\subsection{Compiler Analytics} 
Given a C description of an algorithm, a report is to be generated, highlighting areas within the
code that may see a potential speed-up based on matching to known GPU performant patterns. In order
to achieve this, a static analysis methodology was invoked. Direct source code analysis is
incredibly difficult, and suffers from a number of drawbacks, including the difficulty to parse C
and (notoriously difficult) C++. Considering the time constraints involved it was decided that the
analysis tool would hook into the Clang Libtooling library, which provides access to the Clang
Abstract Syntax Tree, giving a semi-syntax invariant canvas from which to identify relevant parallel
patterns.

Given the difficulty of setting up a generic build environment for the Clang tooling, I decided to
fork a project named OCLint\cite{oclint}, a C, C++ and Obj-C static analysis tool that I had worked
with earlier.  This tool provides a light framework around the Clang tooling, however most
importantly it provides sophisticated build scripts which work on a variety of operating systems and
distributions.

\subsubsection{OCLint Modification} 
Forking the OCLint project, which is licensed under a modified BSD license has saved significant
time and effort from being wasted developing a generic build system around the Clang tooling. The
OCLint software provides a direct method for interacting with the Clang AST by exposing the Clang
Libtooling headers. This increased flexibility has in turn allowed for more time to be spent
developing methods to identify parallel patterns within the generated AST. All changes to the OCLint
software are unrelated to its original intention and design, and as such no pull requests were
lodged, and no modifications I have written have moved upstream. As I have substantially
re-engineered and re-purposed the OCLint software I have elected to give it a new name, the C
Algorithm Parallelisation Analyser (CAPA).

\subsection{GPU Benchmarking}
In order to best provide theoretical performance improvements of algorithms within a codebase, an
analysis of the current hardware available is significantly important. As such rather than just
provide purely theoretical numbers, part of this project involved developing a simple set of GPU
benchmarks which seek to show performance metrics for the identified patterns within the code
analyser. This in effect means that reports generated by the analytics tool may contain specific
information pertaining to the hardware available on the current build and test system. In order to
achieve the best outcome, CUDA was decided on as the framework for development.

\subsubsection{Benchmarks}
GPUs are exceptionally good at high throughput calculations, one particular example is SIMD, meaning
\emph{Same Instruction Multiple Data}. The performance of GPUs and the algorithms they are
particularly useful for is well under continual research, however general problem classes that GPUs
are able to solve efficiently are well understood. These problem sets include algorithms that can be
described by any of the following:

\begin{itemize} 
    \item Map 
    \item Fold/Reduce 
    \item Scan/Prefix Network 
    \item Matrix Operations
\end{itemize}

The actual speed improvements derived from redeveloping serial code to take advantage of the
massively parallel compute power of a GPU differs between each of these operations, however many
serial algorithms have equivalent or more performant alternate parallel implementations. As such
this project involves developing a small set of benchmarks for GPUs that determine their
performance in each of these categories. In order to satisfy time constraints and recognise real
world concerns, I elected to use existing optimised libraries for the individual components of the
benchmarking. I relied on the Eigen library\cite{eigen} for host side matrix operations. On the
device side I relied on a combination of the Thrust template library\cite{thrust} in combination
with CuBLAS\cite{cublas}

\subsubsection{CUDA} 
CUDA is Nvidia's proprietary library and toolchain for developing parallel software. There are 2
main frameworks in the GPU programming space, CUDA and OpenCL. Whilst OpenCL is a FOSS platform, the
development tools are severely lacking in comparison to the CUDA toolkit, and as such it was an easy
decision to follow through with the CUDA. This however has limited the performance metrics to only
comparisons involving CUDA enabled graphics cards. This is not too great a concern however, as the
benchmarking module is highly extensible and as a result can easily integrate with a variety of
backends, including OpenCL or OpenMP.

\pagebreak

%!TEX root = main.tex
%----------------------------------------------------------------------------------------
%	RE-EVALUATION OF INITIAL GOALS
%----------------------------------------------------------------------------------------

\section{Re-evaluation of Initial Goals}
In order to consolidate the current position of this project, and to best identify a pathway to completion, it's important to take another look at the initial requirements set forth in the requirements analysis. Within the requirements analysis there were a set of goals that defined the project and from which all development so far has stemmed; by looking at these requirements and evaluating the current trajectory of the project a detailed description of what is required and how it will be achieved can be compiled.

The requirements analysis broke the project down into 3 major components:

\begin{itemize}
\item GPU Benchmark Development
\item Algorithm Analytics Development
\item Optimisation Analytics Development
\end{itemize}

which then allows the breakdown of what so far has happened in the project.

\subsection{GPU Benchmark Development}
As described earlier, the importance of developing working GPU benchmarking code for known problem classes allows for better analytics and reporting in the serial algorithm analysis portions. This therefore is a key aspect of satisfactorily completing the project. The GPU benchmark development has a number of requirements that describe what the project necessitates.

\subsubsection{Requirements}

\paragraph{[FR.003]} \label{[FR.003]}
\textbf{The program shall run developed benchmark algorithms to further analytical information.}

This requirement relates directly to the overall aim of the project, which is described in the requirements just proceeding this. As the project currently stands there is only one working benchmark developed, it is a best case memory bandwidth test which requests on device memory, fills it with junk host memory, and requests theen now junk device memory be copied back to the host. This test aims to determine the peak memory bandwidth for the device, which can then be used to determine absolute optimal performance of any subsequent alogrithm.

In order to complete the project more benchmarks need to be written, to specifically cover algorithm optimisation cases identified in the serial analytics portion of the project. Necessary benchmarks are as follows: 

\begin{itemize}
\item Peak Memory Bandwidth
\item Peak Map
\item Peak Fold/Reduce
\item Peak Scan
\item Peak Matrix Multiplication
\item Peak Depth first Graph Traversal
\end{itemize}

Upon completion of these benchmarks this requirement will have been completed, interfacing and using the results of these benchmarks are a seperate requirement.

\paragraph{[OA.001]} \label{[OA.001]}
\textbf{The program shall run custom benchmark algorithms to identify GPU
performance.}

This requirement describes that the benchmarking algorithms must be utilised to identify GPU performance. In order to satisfy this requirement my intention is to produce benchmarking code for GPU performance in problem sets that are both known to be performant on a GPU as well as benchmarks that may naively appear to be performant, yet further inspection demonstrates that they are not in fact performant. This is a rather large task in and of itself, and has the potential to be an entire FYP on its own, as such significant compromises must be undertaken. In this particular case only 2 non-performant algorithms will be developed, they are:

\begin{itemize}
\item Recursive Dependent Matrix Calculations
\item Highly Divergent Hashmap Traversal 
\end{itemize}

These problem sets seem to be readily parallelizable, however they in fact exacerbate the weaknesses of the general NVidia GPU architecture (and potentially other co-processor architectures I have not worked with). The results of these benchmarks contextualise the information derived from the code analytics portion, giving rise to more useful metrics defining best and worst case performance.

\paragraph{[OA.002]}
\textbf{The program shall work on all CUDA devices.}

After careful consideration this requirement has been relaxed as it is far too strict. When writing the requirements analysis I was not as familiar with the CUDA toolchain as I am at this point of my project, and as such it is now apparent that writing Compute Capability agnostic code is a very difficult feat. In order to best satisfy the other requirements of this project I shall be limiting benchmarking code to work on CUDA capable devices of Compute Capability 3.5 and above. This compute capability was chosen as the capabilities of CUDA Cards differ significantly pre and post Compute Capability 3. This revision of the initial requirement shall save significant time and effort from being wasted in localisation and highly technical activities that benefit the overall project very little.

\paragraph{[OA.003]}
\textbf{The program shall provide comparative CPU performance metrics.}

This requirement remains as it was initially written, there is no reason why the project should not continue to include a comparison between parallel GPU benchmarks and the relevant serial CPU implementation. This information will be used within the analytics framework.

\paragraph{[OA.004]}
\textbf{The program shall provide a number of different problem class
benchmark algorithms.}

This requirement was covered and expanded under sections \ref{[FR.003]} and \ref{[OA.001]}














\pagebreak

%!TEX root = ./main.tex
%----------------------------------------------------------------------------------------
%	PROJECT STATUS
%----------------------------------------------------------------------------------------

\section{Project Status}

This section provides a summary of what has been completed, and not completed so far within this FYP.


\subsection{Completed Tasks}

\begin{itemize}
\item Stable package and distribution setup
\item Successful identification of simple parallel patterns
\item Initial scaffolding of GPU benchmarking code
\item Basic GPU peak memory performance benchmark
\item Scaffolding of performance metric reporting
\item Capable of analysing arbitrarily large C codebases
\end{itemize}


\subsection{Incomplete Tasks}

\begin{itemize}
\item Identification of more complex parallel patterns
\item Identification of complex parallel patterns
\item Design and coding of GPU benchmarking suite
\item Implementation of skeleton analyser reporter
\item Implementation of GPU performance analyser
\item Implementation of theoretical performance analyser
\item Completion of Poster
\item Completion of presentation video
\item Generation of documentation
\end{itemize}



\pagebreak

%!TEX root = ./main.tex
%----------------------------------------------------------------------------------------
%       CURRENT POSITION REVIEW	
%----------------------------------------------------------------------------------------

\section{Current Position Review}

This section aims to look forward at what needs to be achieved and how what has been done so far sets up the rest of the project for completion.

\subsection{Incomplete Tasks}
\subsubsection{Identification of more complex parallel patterns} 
At the current point within the project, basic parellel patterns are recognisable within given test code. These patterns are Map operations, Reductions and Scan's. Intermediate complexity matrix operations are considered to be a selection of 2D matrix operations. The most optimisable of these is matrix addition, and subtraction, however multiplication is a readily parallelizable task which must be recognised. The current strategy of identifying parallel patterns by utilising the Clang libtooling matcher library has sufficient flexibility to identify these more complex patterns. Given the success of identifying the basic parallel algorithms, it is expected that this will be achieveable within the timeframe set for the project.

\subsubsection{Identification of complex parallel patterns}
Complex parallel patterns primarily refers to identification of breadth first graph traversal, however it also is in reference to identification of parallel inefficient algorithms and recommending a more efficient parallel implementation, such as depth first search, or certain sorting algorithms. These are significantly more difficult to identify given the large variety of ways in which these algorithms may be written. This introduces significant overhead to both the matcher and the callback, and requires significant effort to develop a generic setup for even one of these problem classes. As such these are considered the most complex patterns for identification, and as such require the largest time investment to bring them up to speed.

\subsubsection{Deisgn and coding of GPU benchmarking suite}
The process has already begun for the production of the GPU benchmarking suite. An object heirachy for generic reporting of statistics has been created, and the initial memory benchmark has been implemented, with performances roughly matching theoretical expectations on the test hardware. The benchmarking suite will include problem classes that the algorithm analyser will be targeting, in order to provide the most relevant information for theoretical performance estimation. As it currently stands, this means the production of basic Map, Reduce and Scan kernels, as well as a basic streaming matrix arithmetic kernels, and potentially an implementation of an efficient breadth first search, if time permits.

\subsubsection{Implementation of theoretical performance analyser}
This is the cumulative goal of the entire FYP, in order to best achieve this goal a combination of the source code analyser and the GPU benchmarker is required; by definition this task is currently incomplete as the two major components that comprise this task are also incomplete. In order to achieve this goal compromises must be made between the depth in whcih both aspects of this project are pursued. In order to provide the best possible outcome it is important that good decision making is undertaken in establishing what aspects of the code analyser and benchmarking suite receive the most attention. The key aim of this task is to provide a report detailing parts of the source code that may see improvements from parallelization. This key aim is enhanced by providing further detail about what sort of potential improvements may be seen in a parallelized implementation. Due to the nature of this relationship it is clear that the code analyser is the \textit{master} and the benchmarks are the \textit{slave} components of the project.

\subsubsection{Presentation Related Items}
These naturally rely on aspects of the project being either completed, or close to completion and as such cannot be completed until far later in the project timeline.

\subsection{Completed Tasks}
\subsubsection{Identification of Simple Parallel Patterns}
This is the initial requirement of the code analyser aspect of this project. Currently simple parallel patterns such as map operations have been successfully completed. This is done utilising the Clang libtooling library, using ASTMatchers and the resulting callback. The current AST Matcher for a map operation is:
\lstinputlisting{./Code/mapMatcher.cpp}
Which identifies the key characteristics of map operations, that being; for loops which increment a loop counter by one with an array assignment which is calculated as a one to one mapping from an input array.

Other matchers for simple parallel patterns are similar in structure to the map matcher, as they are all operations over 1 dimensional arrays. Higher dimension data structures require a different matcher structure to account for their increased complexity.

\subsubsection{Inital Scaffolding of GPU Benchmarking Code}
The current scaffolding of the GPU Benchmarking code provides a benchmarking class which defines the interface by which other aspects of this project must interface with the benchmarking code. Currently the interface is:
\lstinputlisting{./Code/benchmarkScaffold.h}
The current interface design provides a clean API for handling all benchmark requests, as well as providing library users the ability to change configuration information between benchmark operations. The key benefit however is within the restricted interaction beyond initial configuration of the benchmark suite, which provides programs utilising the API a simpler interface from which to work. Internal code within the library is up for change, however as the external API aims to remain constant this should not impact potential library users.

Additionally in order to increase code re-use I've attempted to implement the concept of Type-Classes into CUDA C++. This allows for the single case of higher order functions such as Map, Reduce and Scan. This implementation utilises template meta-programming and templated type alias's to produce type safe higher order polymorphism within CUDA compute kernels.
\lstinputlisting{./Code/typeclass.cpp}
This code segment demonstrates typeclass instances for operations over both \textit{Functors} and \textit{BiFunctors}, the \textit{Functor} typeclass is the alias uCat, a function which accepts a single input of type T and returns a single output of type T. The second typeclass definition is the \textit{mCat} defition, defining the monoidal typeclass requirement of \textit{mConCat}. This is any function that accepts 2 inputs of type T and returns an output of type T. These typeclasses are then used to validate the parametric polymorphism of the \textit{biMapKernel}, which implements a polymorphic bimap operation, the key aspect of the \textit{bifunctor} typeclass, as well as the common map kernel, which is the single requirement of the \textit{functor} typeclass. This parametric polymorphism provides an easy to use, clear and concise framework from which to build larger GPU benchmarking routines, by utilising the higher order nature of the GPU \textit{primatives}. To summarise the category theory, this essentially provides a type safe way to write less code.

\subsubsection{Basic GPU Peak Memory Performance Benchmark}
The current peak memory performance benchmark is an incredibly simple memory test, the base benchmark is merely allocates memory on the device, copies memory from the host to the device, performs a no-op on the device and copies the memory back from the device to the host. This is all timed in order to calculate a peak theoretical memory bandwidth. The benchmark was written before the current version of the benchmark object existed, and as such it does not have any dynamic reponse to changing structure. It is merely a proof of concept for the peak memory performance benchmark.
\lstinputlisting{./Code/memoryBandwidth.cpp}
Even though this code does not interface with the GPU benchmark object it is a trivial transformation to effect that change and bring this inital benchmark into a useable form.

\subsubsection{Scaffolding of Performance Metric Reporting}
By utilising the OCLint tool, the rule and reporter interface has already been written. There are a few alterations that still need to be undertaken, as the nature of the reports is somewhat different between the original intention of OCLint, and what I am doing, however the overall scaffolding and relationship between analysis and report is finished, and quite robust.

\subsubsection{Capable of Analysing Arbitrarily Large C Codebases}
As the static code analyser uses the Clang libtooling, if the codebase being run on is capable of being built by clang, then it is possible for the codebase to be analysed CAPA. This is due to the heavy lifting being done by the clang compiler, and the tool hooking in only at the AST stage per compilation unit.



\pagebreak

\end{document}




